\section*{Bab \ref{ch:g2:what-animals-eat} --- Apa yang Dimakan Hewan}

\begin{solution}{exo:g2:what-animals-eat:1}
\omterm{def:g2:what-animals-eat:kinds}{Herbivora}; \omterm{def:g2:what-animals-eat:kinds}{karnivora}; \omterm{def:g2:what-animals-eat:kinds}{omnivora}.
\end{solution}

\begin{solution}{exo:g2:what-animals-eat:2}
(a) \omterm{def:g2:what-animals-eat:kinds}{Herbivora}. (b) \omterm{def:g2:what-animals-eat:kinds}{Karnivora}. (c) \omterm{def:g2:what-animals-eat:kinds}{Omnivora}. (d) \omterm{def:g2:what-animals-eat:kinds}{Herbivora}. (e)
\omterm{def:g2:what-animals-eat:kinds}{Omnivora}.
\end{solution}

\begin{solution}{exo:g2:what-animals-eat:3}
Penggilas yang rata --- \omterm{def:g1:plants-around-us:plant}{tumbuhan} seperti rumput itu liat dan harus
digilas lama sebelum ditelan; tidak ada mangsa yang perlu diterkam.
\end{solution}

\begin{solution}{exo:g2:what-animals-eat:4}
Pendek dan tebal: memecah \omterm{def:g2:from-seed-to-plant:seed}{biji}. Melengkung dan berkait: merobek
daging. Panjang dan tipis: mencungkil serangga dari kulit kayu dan
celah-celah.
\end{solution}

\begin{solution}{exo:g2:what-animals-eat:5}
Kutu \omterm{def:g1:plants-around-us:parts}{daun} yang mungil --- kepik memakannya dalam jumlah besar, dan
itulah sebabnya tukang kebun menyukainya.
\end{solution}

\begin{solution}{exo:g2:what-animals-eat:6}
Rumput memberi sedikit zat makanan pada tiap suapan, jadi sapi harus
merumput berjam-jam untuk mendapat cukup; daging itu kaya, jadi satu
santapan yang baik dapat menahan rubah cukup lama.
\end{solution}

\begin{solution}{exo:g2:what-animals-eat:7}
Cacing dan serangga pada musim panas; buah beri pada musim dingin,
ketika tanahnya terlalu keras bagi cacing. Daftar makanan yang lebar
memungkinkan seorang \omterm{def:g2:what-animals-eat:kinds}{omnivora} berganti makanan ketika yang satu habis
--- pertolongan besar pada musim yang sulit.
\end{solution}

\begin{solution}{exo:g2:what-animals-eat:8}
Jawaban terbuka. Biasanya: burung pipit dan burung finch (paruh pendek
tebal) mengambil \omterm{def:g2:from-seed-to-plant:seed}{biji}; burung gelatik mencungkil serangga dan \omterm{def:g2:from-seed-to-plant:seed}{biji};
burung hitam menusuk tanah mencari cacing dan mengambil buah beri.
Kecocokan paruh dengan daftar makanan pada proposisinya berlaku dengan
baik.
\end{solution}

\begin{solution}{exo:g2:what-animals-eat:9}
Daftar makanan seorang \omterm{def:g2:what-animals-eat:kinds}{karnivora}: \omterm{def:g1:animals-around-us:animal}{hewan} lain. Langkah 1 dan 2 ---
giginya (penerkam yang runcing) dan perkakas berburunya (cakar, \omterm{def:g1:the-five-senses:sense}{mata}
yang menghadap ke depan); langkah 4 menyebutkan namanya.
\end{solution}

\begin{solution}{exo:g2:what-animals-eat:10}
Rubah memakan kelinci, dan kelinci memakan rumput. Tanpa rumput, tanpa
kelinci; tanpa kelinci, rubah kelaparan. Rubahnya bergantung pada
rumput \emph{melalui} mangsanya --- santapan di alam membentuk rantai.
\end{solution}
